\section{Proving Computable Enumerability}

The final theorem requires more than abstract monotone limits.  The sets
\texttt{Aset} and \texttt{Bset} must satisfy the repository definition
\texttt{CE}, i.e. each must be the halting domain of an oracle-free local
program.  The file \texttt{CE.lean} proves this by combining a
computability-friendly encoded construction with the Mathlib bridge.

\subsection{Encoded Construction}

The mathematical construction state contains custom structures
\texttt{ReqState} and \texttt{ConsState}.  Rather than building custom
\texttt{Primcodable} instances for these types, \texttt{CE.lean} defines a
parallel tuple representation:

\begin{lstlisting}
abbrev RS : Type := Option Nat * Bool * Nat
abbrev St : Type := List Nat * List Nat * List RS * Nat

def encReq (r : ReqState) : RS :=
  (r.witness, r.acted, r.restraint)

def encSt (st : ConsState) : St :=
  (st.Alist, st.Blist, st.reqs.map encReq, st.fresh)
\end{lstlisting}

It then defines tuple-level versions of the convergence check, attention
predicate, idle step, appointment step, action step, stage step, and full
stage recursion: \texttt{convCheckN}, \texttt{reqAttN}, \texttt{idleSt},
\texttt{appointSt}, \texttt{actSt}, \texttt{stepN}, and \texttt{stageN}.

\subsection{Simulation}

The key simulation theorem is:

\begin{lstlisting}
theorem encSt_stepState (st : ConsState) (s : Nat) :
    encSt (stepState st s) = stepN (encSt st) s
\end{lstlisting}

By induction, \texttt{encSt_stage} proves that the encoded construction
computes the encoding of the mathematical construction at every stage.
The lemmas \texttt{mem_Alist_iff_stageN} and
\texttt{mem_Blist_iff_stageN} transfer membership facts between the two
representations.

Two list-surgery lemmas, \texttt{map_range_set} and \texttt{map_range_act},
explain how list update and lower-priority initialization are represented
as maps over \texttt{List.range}.  This representation makes the operations
amenable to Mathlib's primitive-recursive list combinators.

\subsection{Primitive Recursiveness and Semidecision}

\texttt{CE.lean} proves primitive recursiveness of the encoded transition
and stage function.  The proof depends on \texttt{nrun_primrec}, the
primitive-recursive evaluator on numbered oracle codes and finite
snapshots.  The chain of results is:

\[
\begin{array}{c}
\text{Mathlib Primrec/Partrec closure lemmas}\\
\downarrow\\
\text{\texttt{nrun\_primrec} and primitive-recursive encoded step}\\
\downarrow\\
\text{\texttt{stageN} is primitive recursive}\\
\downarrow\\
\text{membership in encoded stage lists is decidable computably}\\
\downarrow\\
\exists s,\ x\in A_s \text{ and } \exists s,\ x\in B_s
\text{ are semidecidable.}
\end{array}
\]

The final c.e. theorems are:

\begin{lstlisting}
theorem ce_Aset : CE Aset
theorem ce_Bset : CE Bset
\end{lstlisting}

Each proof builds a Mathlib partial-recursive search over stages using
\texttt{Partrec.rfind}, then applies \texttt{ce_of_partrec_dom} from
\texttt{MathlibBridge.lean} to obtain the repository-local \texttt{CE}
statement.

